\honest{Feeling Moral}{Eliezer Yudkowsky}{2015}

Something is killing people, and you can save 400 lives for certain, or save 500 with probability 90\% and none otherwise. ``Most people choose option 1,'' and I call that foolish, since the gamble saves 450 lives on average, and lives saved do not diminish in marginal value. ``How can you gamble with human lives?'' you cry. What if the 10 per cent strikes? ``You're following your rationality off a cliff!'' \nb{The experiment behind this is Tversky and Kahneman's ``Asian disease'' problem (1981), in which 72 per cent chose to save 200 of 600 people for certain. There the two options had the same expected value. The numbers that make the sure choice foolish are the post's own, and no study is cited for them.}

You may cry that I gamble with human lives. But describe the same options as deaths, 100 dead for certain against a 10\% chance that all 500 die, and a majority choose the gamble, ``Even though it's the same gamble.'' A sure gain feels better than an unsure one, and a sure loss feels worse. And one can grandstand on that framing too: how can you condemn 100 people to certain death? Everyone makes it, or no one does! \nb{This reversal is Tversky and Kahneman's central finding, and it makes the post's point whichever option is better: when the choice flips with the wording, at least one of the choices is following the words.}

A life, with its joys and pains over decades, is worth far more than your comfort with a plan. If computing expected utility feels cold-blooded, that feeling is not even a feather in the scales: ``Just shut up and multiply.''

Now take a googolplex, 10 to the power of a googol, which is itself $10^{100}$: a googolplex people with hiccups, against one person slowly torn apart by sadistic mutant sharks. If a hiccup has ``any'' negative value, then ``on pain of decision-theoretic incoherence'' some number of hiccups must be worse than the shark attack. \nb{This is Yudkowsky's ``Torture vs.\ Dust Specks'' question in new words; that post is not in the book. Coherence alone does not force the conclusion: a ranking in which each hiccup makes things worse, but the total never reaches the shark attack, is consistent. The conclusion needs the premise that each hiccup counts the same however many there are, which the next post argues for. The first version of this post argued with a chain of small steps from torture to dust specks; this version leaves the chain out.}

Such dilemmas are not blood sports for philosophers at dinner parties but everyday life: \$50 on a console game or to charity, one bone marrow transplant for \$700,000 or bed nets that prevent the deaths of some 200 children. Research shows that people call lives ``sacred'' and are indignant at any trade against money, sometimes wanting to punish whoever proposed it. Many take pride in looking away from such trade-offs. \nb{The research is Philip Tetlock's. Tetlock also reports that people often accept such trades when they are described as routine or tragic.}

My favourite anecdote: an agency rejected a report that priced a life-saving project, because you cannot put a dollar value on a life, and then decided not to carry out the project. Such trades feel awful. \nb{No source is given. In the first version Yudkowsky wrote: ``my books are packed at the moment, so no citation for now.''}

Altruism is not a warm feeling; doing it for the spiritual benefit ``is nothing but selfishness.'' The primary thing is to help others, whatever the means. If there seems a fierceness in this maximizing, a small cold flame at the centre of rationality, the other way might feel better inside you, ``But it wouldn't work.'' If you give yourself to rationality without holding back, ``rationality gives to you in return.'' But that works only if you do not tell yourself you would feel better being less rational. Should you be sad to have the chance to actually help people? You cannot reach your potential if you regard your gift as a burden.
